\subsubsection{Actividades [Pendientes: 1. confirmar participantes y tareas de evaluación]}\label{sec:transicion-actividades}
Se prepararán tareas representativas del flujo existente y se definirán las condiciones de aplicación, los resultados esperados y la forma de registrar observaciones. Durante las sesiones se examinará si los participantes logran plantear una consulta, reconocer la información espacial y temporal presentada e identificar los límites de la estimación, evitando solicitar acciones que excedan las funciones comprometidas. Los hallazgos se incorporarán al informe E6 y se tratarán mediante correcciones verificables de la versión evaluada. Después de cada ajuste pertinente se repetirán las comprobaciones afectadas, conservando la relación entre problema observado, modificación y resultado de la nueva prueba.

Se consolidarán el código y el ejecutable E5, las configuraciones E3 y el catálogo E4 necesarios para reproducir la demostración, y se revisará la correspondencia final con E1 y E2. El manual E7 describirá instalación, dependencias, uso del flujo, condiciones representadas y límites conocidos, mientras que el informe final reunirá resultados y conclusiones sustentados por las evidencias obtenidas. La entrega al director incluirá una demostración preparada con datos y versiones identificadas, de manera que pueda distinguirse lo ejecutado de las funciones o condiciones que hayan quedado fuera de la evaluación.
